% Mathématiques 4e — cours V3
% Probabilités et simulation

\section{Objectifs}

\begin{itemize}
\item Utiliser le vocabulaire : expérience aléatoire, issue, événement.
\item Déterminer des probabilités dans des situations simples, équiprobables ou non.
\item Calculer la probabilité d'un événement à partir d'un modèle explicite.
\item Comprendre et utiliser l'événement contraire.
\item Vérifier que les probabilités de toutes les issues forment un total égal à $1$.
\item Répéter ou simuler une expérience.
\item Calculer des fréquences expérimentales.
\item Observer la stabilisation des fréquences lorsque le nombre d'essais augmente.
\item Étudier des expériences très simples à deux étapes avec un tableau à double entrée.
\end{itemize}

\section{1. Expérience aléatoire}

Une expérience aléatoire possède plusieurs résultats possibles sans que l'on puisse prévoir avec certitude lequel se réalisera.

Chaque résultat élémentaire est une issue.

Un événement regroupe une ou plusieurs issues.

Exemple avec un dé :
\[
A=\text{« obtenir un nombre supérieur à }4\text{ »}.
\]

Les issues favorables sont :
\[
5,\ 6.
\]

\section{2. Probabilité}

À chaque événement $A$, un modèle associe un nombre :
\[
P(A)
\]
compris entre :
\[
0
\]
et :
\[
1.
\]

\[
\boxed{0\le P(A)\le1}.
\]

Une probabilité n'est pas une prédiction certaine du prochain résultat.

\section{3. Équiprobabilité}

Si toutes les issues sont également probables :
\[
P(A)=\dfrac{\text{nombre d'issues favorables}}{\text{nombre total d'issues}}.
\]

Avec un dé équilibré :
\[
P(\text{pair})=\dfrac36=\dfrac12.
\]

\section{4. Situation non équiprobable}

Une roue comporte trois secteurs dont les angles valent :
\[
180^\circ,\quad120^\circ,\quad60^\circ.
\]

Si le modèle suppose une probabilité proportionnelle à l'angle du secteur :
\[
P(A)=\dfrac{180}{360}=\dfrac12,
\]
\[
P(B)=\dfrac{120}{360}=\dfrac13,
\]
\[
P(C)=\dfrac{60}{360}=\dfrac16.
\]

Il y a trois issues, mais leurs probabilités ne sont pas égales.

\begin{attention}
Le nombre d'issues ne suffit pas à justifier une équiprobabilité.
\end{attention}

\section{5. Somme des probabilités}

Pour des issues qui couvrent tous les résultats possibles :
\[
\boxed{p_1+p_2+\cdots+p_n=1}.
\]

Dans l'exemple de la roue :
\[
\dfrac12+\dfrac13+\dfrac16=1.
\]

Cette égalité permet de contrôler un modèle.

\section{6. Événement contraire}

L'événement contraire de $A$ se réalise exactement lorsque $A$ ne se réalise pas.

\coursefigure{/images/mathematiques/4eme/probabilites-contraire.svg}{Représentation d'un événement A et de son événement contraire partageant l'ensemble des issues.}{Un événement et son contraire sont incompatibles et couvrent ensemble tous les résultats possibles.}

\begin{propriete}
\[
\boxed{P(\overline A)=1-P(A)}.
\]
\end{propriete}

\section{7. Exemple de contraire}

Avec un dé équilibré :
\[
A=\text{« obtenir au moins }5\text{ »}.
\]

\[
P(A)=\dfrac26=\dfrac13.
\]

L'événement contraire est :
\[
\overline A=\text{« obtenir un nombre inférieur ou égal à }4\text{ »}.
\]

Donc :
\[
P(\overline A)=1-\dfrac13=\dfrac23.
\]

\section{8. Événements incompatibles}

Deux événements sont incompatibles s'ils ne peuvent pas se produire en même temps.

Avec un seul lancer de dé :
\[
A=\text{« obtenir }2\text{ »}
\]
et :
\[
B=\text{« obtenir }5\text{ »}
\]
sont incompatibles.

Dans ce cas :
\[
P(A\text{ ou }B)=P(A)+P(B).
\]

\section{9. Fréquence expérimentale}

Si un événement se produit :
\[
n_A
\]
fois lors de :
\[
n
\]
essais, sa fréquence est :
\[
\boxed{f_A=\dfrac{n_A}{n}}.
\]

Par exemple :
\[
43
\]
succès sur :
\[
100
\]
essais donnent :
\[
f=0{,}43=43\%.
\]

\section{10. Probabilité et fréquence}

La probabilité appartient au modèle.

La fréquence appartient aux observations.

Une pièce équilibrée possède :
\[
P(\text{pile})=0{,}5.
\]

Une série de $20$ lancers peut pourtant donner :
\[
f(\text{pile})=0{,}65.
\]

Ce résultat n'est pas contradictoire avec le modèle.

\section{11. Stabilisation des fréquences}

\coursefigure{/images/mathematiques/4eme/probabilites-stabilisation.svg}{Courbe de fréquence se rapprochant progressivement d'une valeur de référence.}{Les fluctuations sont fortes au début ; lorsque le nombre d'essais augmente, la fréquence tend à devenir plus stable autour de la probabilité du modèle.}

Cette observation peut être étudiée matériellement ou par simulation.

\section{12. Simulation}

Un tableur ou un programme permet de simuler rapidement :
\[
100,\ 1000,\ 10\,000
\]
expériences.

On peut enregistrer à chaque étape :
\begin{itemize}
\item l'effectif d'un événement ;
\item sa fréquence ;
\item l'écart à la probabilité théorique.
\end{itemize}

La simulation permet d'observer un comportement statistique.

\section{13. Pseudo-code simple}

\begin{verbatim}
succès = 0
répéter 1000 fois
    simuler une issue
    si l'événement A se réalise
        succès = succès + 1
fréquence = succès / 1000
\end{verbatim}

La condition sert à compter les réalisations de l'événement.

\section{14. Modèle à probabilités données}

Une urne contient des objets choisis selon un mécanisme qui donne :
\[
P(\text{rouge})=0{,}5,
\]
\[
P(\text{bleu})=0{,}3,
\]
\[
P(\text{vert})=0{,}2.
\]

On n'a pas besoin de supposer l'équiprobabilité : les probabilités sont fournies par le modèle.

\section{15. Retrouver une probabilité manquante}

Si quatre issues ont pour probabilités :
\[
0{,}25,\quad0{,}15,\quad0{,}40,\quad p,
\]
alors :
\[
0{,}25+0{,}15+0{,}40+p=1.
\]

Donc :
\[
p=0{,}20.
\]

\section{16. Deux étapes très simples}

On peut étudier une expérience comprenant deux étapes lorsque le nombre de cas reste faible.

Exemple :
\begin{itemize}
\item lancer une pièce ;
\item choisir ensuite un nombre parmi $1$ et $2$.
\end{itemize}

Les résultats peuvent être organisés dans un tableau à double entrée.

\section{17. Tableau à double entrée}

\begin{tabular}{c|c|c}
 & $1$ & $2$ \\
Pile & $(P,1)$ & $(P,2)$ \\
Face & $(F,1)$ & $(F,2)$ \\
\end{tabular}

Si les quatre couples sont équiprobables, chacun possède une probabilité :
\[
\dfrac14.
\]

\begin{remarque}
En 4e, on privilégie les situations très simples. L'arbre pondéré de probabilités n'est pas un outil central attendu ici.
\end{remarque}

\section{18. Événement sur deux étapes}

Dans le tableau précédent :
\[
A=\text{« obtenir pile ou le nombre }2\text{ »}.
\]

Les couples favorables sont :
\[
(P,1),\quad(P,2),\quad(F,2).
\]

Donc :
\[
P(A)=\dfrac34
\]
dans le modèle équiprobable.

\section{19. Expérience réelle et modèle}

Une simulation ne prouve pas qu'un modèle est exact.

Si une pièce supposée équilibrée donne :
\[
560
\]
piles sur :
\[
1000
\]
lancers, la fréquence :
\[
0{,}56
\]
est une observation.

L'interprétation dépend du contexte et du nombre d'essais.

\section{20. Lire une affirmation probabiliste}

Dire :
\[
P(A)=0{,}8
\]
signifie que le modèle attribue une forte probabilité à $A$.

Cela ne signifie pas que, sur exactement $10$ essais, $A$ se réalisera obligatoirement :
\[
8
\]
fois.

Les fréquences fluctuent.

\section{21. Contrôler un résultat}

Une probabilité :
\[
1{,}2
\]
est impossible.

De même :
\[
-0{,}1
\]
est impossible.

Dans un modèle complet, la somme des probabilités des issues doit valoir :
\[
1.
\]

\section{22. Méthode}

\begin{methode}
Pour résoudre un problème :
\begin{enumerate}
\item identifier l'expérience et les issues ;
\item déterminer si le modèle est équiprobable ou si des probabilités sont données ;
\item définir précisément l'événement ;
\item calculer sa probabilité ;
\item utiliser l'événement contraire si cela simplifie ;
\item contrôler que le résultat appartient à $[0;1]$.
\end{enumerate}
\end{methode}

\section{23. Erreurs fréquentes}

\begin{itemize}
\item Supposer l'équiprobabilité parce qu'il y a plusieurs issues.
\item Confondre fréquence et probabilité.
\item Oublier que les probabilités des issues totalisent $1$.
\item Calculer le contraire par $-P(A)$ au lieu de $1-P(A)$.
\item Additionner des probabilités d'événements qui se recouvrent sans analyser les cas.
\item Croire qu'une simulation garantit exactement la fréquence théorique.
\item Construire un outil complexe alors qu'un tableau suffit.
\end{itemize}

\section{À retenir}

\begin{aretenir}
Toute probabilité est comprise entre $0$ et $1$.

Pour l'événement contraire :
\[
\boxed{P(\overline A)=1-P(A)}.
\]

Une fréquence expérimentale vaut :
\[
\boxed{f=\dfrac{\text{nombre de réalisations}}{\text{nombre d'essais}}}.
\]

Lorsque le nombre d'essais augmente, les fréquences tendent à se stabiliser autour des probabilités du modèle.

L'équiprobabilité doit être justifiée et non supposée.
\end{aretenir}
